\chapter{Einleitung}
\label{ch:einleitung}

\section{Problemstellung}
\label{sec:motivation}
Die vorausschauende Instandhaltung (engl. \textit{Predictive Maintenance}; kurz \textit{PdM}) bietet für viele Instandhaltungsarbeiten großes Verbesserungspotential. Insbesondere im Zusammenhang mit der Industrie~{4.0} hat PdM in den letzten Jahren viel an Bedeutung gewonnen. PdM rückt für zunehmend mehr Anwendungen in den Fokus, ist aber auch noch heute in vergleichsweise wenigen Anwendungen zu finden. 

Ein Grund dafür warum PdM weniger verbreitet ist als z.B. präventive Methoden ist zum einen der hohe Investitionsaufwand und zum anderen die Tatsache, dass in vielen Bereichen mögliche Anwendungsfälle nicht bekannt sind. Zwar ist es denkbar, dass für den speziellen Fall eine andere Herangehensweise die Bessere ist, jedoch selbst um dies festhalten zu können, muss das Potential für die jeweilige Anwendung analysiert werden. 

Grundsätzlich bedarf es einer umfassenden Analyse eines neuen PdM-Usecases ehe dieser Umgesetzt werden sollte.

Die vorausschauende Instandhaltung bietet im Vergleich zur präventiven Instandhaltung das Potentiale der Senkung der Instandhaltungskosten und einen darüber hinausgehenden Informationsgewinn der auch für Prozesssteuerung und Verbesserungen genutzt werden kann. Dies ist aber nur unter bestimmten Voraussetzungen der Fall, welche sich dazu vom betrachteten Fall abhängen.

Um eine Einschätzung über das Potential eines bestimmten Anwendungsfall abgeben zu können, ist es notwendig diesen so weit zu erarbeiteten, dass ein Vergleich mit dem zu ersetzenden Verfahren sinnvoll durchgeführt werden kann.

Für die vorausschauende Wartung sind Vorhersagemodelle nötig. Deren Qualität ist einer der wichtigste Faktor für den Nutzen des gesamten Anwendungsfalls. Ist beispielsweise die Verhinderung von Ausfällen und eine besserte Ausnutzung der Bauteillebensdauer das Ziel, so muss das Model Vorhersagen treffen die zu einem gewissen Maß zutreffend sind, um einem anderen Verfahren gegenüber bevorzugt werden zu können.

Mit dem Aufschwung der Industrie 4.0 und der damit einhergehenden stetig größer werdenden Menge verfügbarer Daten, rückt PdM weiter in den Fokus. Schließlich bieten sich durch eine große Menge Daten neue Möglichkeiten. Nicht nur können bestehende Anwendungen mit größerer Präzision ausgeführt werden; es ergeben sich auch grundlegend neue Anwendungsmöglichkeiten.

PDM ist in vielen technischen Bereichen insbesondere deshalb interessant, weil  dadurch folgende Potentiale erschlossen werden können \cite[S.~31]{Holzel.2019}. Diese sind in Abgrenzung zur präventiven Instandhaltung zu sehen:
\begin{itemize}
	\item Bedarfsgerechte Ersatzteilbereitstellung
	\item Instandhaltungsmaßnahmen können zum optimalen Zeitpunkt durchgeführt werden
	\item Verlängerung der Anlagen- und Bauteillebensdauer
	\item Reduzierte Stillstandzeiten
	\item Reduzierte Instandhaltungskosten
	\item Verbesserung der Anlagenzuverlässigkeit
\end{itemize}

Auch für die Instandhaltung von Schienenfahrzeugen erscheinen PDM-Anwendungen vielversprechend. Neben den oben genannten Punkten erscheint durch die Verringerung von Ausfällen auch Imageschaden abwendbar. Dieser Entsteht, wenn  der Ausfall einer Anlage (z.B. Zug oder Tür) sich direkt auf den Kunden auswirkt; in diesem Fall der Fahrgast.

Ein für Fahrgäste zu beobachtender Fall sind Türstörungen. In der Regel äußern sich diese darin, dass die Zugtür nicht mehr entsprechend der Sicherheitsanforderungen nutzbar sind. In den meisten Störungsfällen genügt es die Tür außer Betrieb zusetzen und um so die Sperrung eines Waggons oder gar den Ausfall einer ganzen Fahrt zu verhindern. Dennoch ist der Ausfall einer Zugtür für den Kunden direkt ersichtlich und führt damit zu einem Imageschaden.

Die Fragestellung, die bei der Umsetzung aller PDM-Anwendungen zunächst geklärt werden muss ist, ob das entwickelte Modell zur Vorhersage ausreichend präzise Vorhersagen trifft. Insbesondere entscheidend ist dabei der vergleich mit der gegenwärtigen Instandhaltungsstrategie.

In jedem Fall muss die Frage geklärt werden, ob sich die Umstellung auf ein vorrauschauendes Instandhaltungsprogramm lohnt. Ein umfangreiche Analyse der zu erwartenden Kosten stellt die Grundlage zur Beantwortung dieser Frage dar. Das Vorhersagemodel, auf dem der Anwendungsfall basiert, stellt die wichtigste Grundlage zur Kostenanalyse dar. Ohne eine zuverlässige Einschätzung über die Treffsicherheit der Vorhersagen, kann kein Kostenvergleich zu einem anderen Verfahren erstellt werden. Die Entwicklung eines Vorhersagemodels stellt deswegen eine entscheidende Grundlage dar. 

% \begin{figure}[ht]
% 	\centering
% 	\includegraphics[scale=0.4]{FH_Kiel_Logo_deut_rgb.jpg}
% 	\caption{Logo}
% 	\label{fig:fhlogo}
% \end{figure}

\section{Ziel der Arbeit}
\label{sec:ziel}
Zunächst soll eine neuer PdM-Usecase beschrieben werden, der auf die Erkennung von beschädigten Bauteilen einer Türschließanlage abzielt. Dabei soll lediglich eine Klassifikation mit zwei Kategorien; \enquote{intakt} und \enquote{beschädigt} möglich sein. In diesem Zusammenhang sollen auch Kriterien beleuchtet werden, an denen eine erfolgreiche Umsetzung des Usecases gemessen werden kann.

Um ein Klassifikationsmodel entwickeln zu können, mit dem der Usecase Umgesetzt werden kann, wird ein Datensatz aufgenommen. Dazu werden an der Türschließanlage Messungen durchgeführt. Für intakte und beschädigte Umlenkrollen wird eine Auswahl an Messwerten aufgenommen. 

Im Zuge des \textit{Dataunderstanding} wird die Bedeutung der Daten für den Usecase diskutiert. Dabei wird außerdem ein Datensatz erarbeitet, der den rohen Datensatz beschreibend zusammenfasst. Ziel ist dabei den Informationsgehalt der Daten zu bündeln und so die Verallgemeinerungsleistung der Modelle zu fördern.

Um ein Model zu bestimmen, das die besten Vorhersagen über den Zustand der Umlenkrollen liefert, wird zunächst eine Vorauswahl infrage kommender Modeltypen durchgeführt. Anschließend werden Modelle der in frage kommenden Typen erstellt. Dabei wird für jedes Model eine Parameterstudie durchgeführt.

Nach der Modellierung erfolgt eine Bewertung der Modelle. Anhand ausgewählter Kriterien, die für den Usecase relevant sind, wird das am besten geeignete Model bestimmt.

Abschließend soll ein Ausblick gegeben werden, der mögliche nachfolgende Schritte beleuchtet, um den Anwendungsfall praxistauglich zu gestalten, sollte ein geeignete Vorhersagemodel entwickelt werden können.

\section{Vorgehensweise}
\label{sec:vorgehensweise}
In den Kapiteln 2 bis 4 wird der Stand der Technik der Elemente beschrieben, die für den Inhalt der Arbeit von Bedeutung sind. \cref{ch:instandhaltungsstrategien} befasst sich mit den grundlegend mit der vorausschauenden Instandhaltung und diskutiert dessen Potentiale und Voraussetzungen.

In \cref{ch:machine_learning} werden zunächst Kernpunkte von maschinellem Lernen erörtert, ehe auf die Grundlagen der Lernmodelle eingegangen wird, die im weiteren Verlauf dieser Arbeit verwendet werden. Für die einzelnen Modelle werden u.a. Funktionsweise, Parametereinstellungen, Vor- und Nachteile, sowie die Eignung für den in \cref{ch:usecase} dargestellten Anwendungsfall besprochen. 

\cref{ch:usecase} stellt den neu erdachten PDM-Usecase vor. Dabei wird an den entsprechenden Stellen auf die Inhalte aus \cref{ch:instandhaltungsstrategien} verwiesen und ein Beispielszenoario beschrieben, wie die Anwendung des Usecases aussehen könnte. Um den Usecase in den Kontext der praktischen Bahntechnik zu setzen, werden die Inhalte zweier Interviews mit Bahnsachverständigen rekapituliert. In Zuge dessen wird auch der reale Bedarf des vorgeschlagenen Usecases diskutiert. Die Darstellung des Usecases schließt mit einer Aufführung der Erfolgskriterien des Usecases.

In \cref{ch:methodik} wird die Methodik beschrieben mit der der Usecase exemplarisch untersucht wurde. Es wird der Versuchaufbau beschrieben, eine Einschätzung über den Datensatz abgegeben und eine begründete Vorauswahl maschineller Lernmodelle getroffen (vgl. \cref{ch:machine_learning}). In \cref{sec:bewertungskriterien} werden die Erfolgskriterien des Usecases aus \cref{ch:usecase} aufgegriffen, um die zu erstellenden Modelle daran bewerten zu können. In \cref{sec:bewertungskriterien} wird damit eine Präferenzfunktion zur Beurteilung der Modelle erstellt.

In \cref{sec:modellierung} erfolgt die Erstellung der maschinellen Lernmodelle anhand des aufgenommen Datensatzes. Um möglichst gute Modelqualitäten zu erreichen wird während des Trainieren der Modelle eine Parameterstudie durchgeführt.

In \cref{ch:modelbewertung} wird die Eignung der trainierten Modelle für den Usecase dargestellt und diskutiert. Es wird eine Einschätzung über die Qualität der Modelle abgegeben und anhand der Präferenzfunktion das Model bestimmt welches am besten für den Usecase geeignet ist.

\cref{ch:fazit} fasst die Ergebnisse der Arbeit zusammen und gibt einen Ausblick über Verbesserungsmöglichkeiten in Bezug auf den Usecase an sich und die erarbeiteten Ergebnisse.